% !TeX root = ../main.tex
\section*{Выводы}
В ходе выполнения лабораторной работы было проведено экспериментальное исследование точности позиционирования навигационной аппаратуры потребителей (НАП) ГНСС в различных городских условиях. На основе полученных данных можно сделать следующие выводы:

\begin{enumerate}
    \item \textbf{Влияние условий окружающей среды:} Подтверждено, что характер застройки (высотность и плотность) оказывает существенное влияние на точность определения местоположения и стабильность вычисленных координат.
    \item \textbf{Открытая местность:} На открытом пространстве (в радиусе 50 метров отсутствуют здания) погрешность составила 14--15 метров при 11 обнаруженных спутниках. Отсутствие препятствий минимизирует переотражения, что делает расчетную точку более стабильной.
    \item \textbf{Влияние малоэтажной застройки:} Вблизи двухэтажного здания абсолютная ошибка составила 13--14 метров (радиус вероятной области 26 метров), однако наблюдались постоянные скачки координат. Это наглядно демонстрирует негативное влияние эффекта многолучевости: приемник смартфона обрабатывал суперпозицию прямого и отраженного от фасада сигналов, что привело к сильным флуктуациям.
    \item \textbf{Влияние многоэтажной застройки:} Вблизи высотного здания (28 этажей) зафиксирована максимальная ошибка позиционирования --- 26 метров. Примечательно, что количество обнаруженных спутников увеличилось до 16, однако сильное затенение значительной части небосвода и мощные переотражения от массивного фасада привели к наихудшей точности навигационного решения среди всех проведенных опытов.
    \item \textbf{Общий итог:} Наличие высотных и отражающих объектов в городских условиях критически снижает точность НАП ГНСС. Эффект многолучевости является основным фактором, приводящим к искажению координат на величины порядка нескольких десятков метров, что необходимо учитывать при использовании навигационных систем в условиях плотной городской застройки.
\end{enumerate}